\textbf{Proof.}
By \texttt{def:zhao-well-defined-ordinary-domain}, \(\widetilde{\mathcal B}_Z\) is precisely Zhao's unquotiented positive-definite rank-two ordinary-inverse restriction, and \(q_Z\) maps it onto \(\mathcal B_Z\).  The complete rank-two part of \texttt{thm:triangle-completion-separation-certificate} gives
\[
\operatorname*{argmax}_{[C]\in\mathcal B_Z}g_C(0)=\mathcal A_0,
\qquad
\operatorname*{argmin}_{[C]\in\mathcal B_Z}w(C)=\mathcal A_{\mathrm{mm}},
\]
with respective optimal values \(31/80\) and \(2613/8432\), and it gives
\[
w(C_0)=\frac{521}{1200}
\qquad\text{for every }C_0\in\mathcal A_0.
\tag{1}
\]

The map \(q_Z\) is surjective by \texttt{thm:zhao-ordinary-domain-quotient-transfer}.  The same theorem proves that \(g_{q_Z(B)}(0)\) and \(w(q_Z(B))\) are constant on each fiber of \(q_Z\).  Therefore, for the gain criterion,
\[
\begin{aligned}
B\in\operatorname*{argmax}_{C\in\widetilde{\mathcal B}_Z}g_{q_Z(C)}(0)
&\Longleftrightarrow
q_Z(B)\in\operatorname*{argmax}_{[C]\in\mathcal B_Z}g_C(0)\\
&\Longleftrightarrow q_Z(B)\in\mathcal A_0,
\end{aligned}
\]
which is exactly
\[
\operatorname*{argmax}_{B\in\widetilde{\mathcal B}_Z}g_{q_Z(B)}(0)
=q_Z^{-1}(\mathcal A_0).
\tag{2}
\]
Applying the identical fiber argument to the loss criterion gives
\[
\operatorname*{argmin}_{B\in\widetilde{\mathcal B}_Z}w(q_Z(B))
=q_Z^{-1}(\mathcal A_{\mathrm{mm}}).
\tag{3}
\]
Surjectivity and fiber invariance also preserve the corresponding objective values.  In particular, (1)--(3) imply that every unquotiented zero-covariance optimizer has worst-endpoint loss \(521/1200\), whereas every unquotiented minimax optimizer has loss \(2613/8432\).  Direct subtraction yields
\[
\frac{521}{1200}-\frac{2613}{8432}
=\frac{4912}{39525}>0.
\tag{4}
\]

Equations (2)--(4) use only the positive-definite rank-two restriction.  Rank-zero and rank-one classes are not elements of \(\widetilde{\mathcal B}_Z\); the separate conclusions about them in the triangle completion theorem are conclusions on Zhao's published all-block-separable pseudoinverse carrier expressed through the note's projective quotient \(\mathcal B_Z^+\).  This proves the corrected domain attribution and all asserted optimizer and value identities. \(\square\)